\scanpage{8}
\mspage{III.7}
of $\Gamma$, then $\Gamma^\Lambda=\mathrm{Hom}_{\Lambda^{\mathrm e}}(\Lambda,\Gamma)$ is called the \uline{commutant of $\Lambda$ in $\Gamma$} and is a subring of $\Gamma$. In [4, Theorem 3.5] the following is shown : if $\Gamma$ is central separable over $\mathrm A$, and $\Lambda$ and $\Omega$ are subalgebras of $\Gamma$ such that $\Lambda\otimes_{\handwritten{A}}\Omega\approx\Gamma$ under the mapping $\mathrm{x}\otimes\mathrm{y}\longrightarrow\mathrm{xy}$, then $\Lambda$ is the commutant of $\Omega$ and $\Omega$ is the commutant of $\Lambda$ in $\Gamma$.

Analogously, we can define the commutant of a subalgebra of sheaves. If $\handwritten{\mathscr{G}}$ is a sheaf of rings with center $\handwritten{\mathscr{A}}$, and $\mathscr{L}$ is an $\handwritten{\mathscr{A}}$-subalgebra of $\handwritten{\mathscr{G}}$, we define the commutant of $\mathscr{L}$ in $\handwritten{\mathscr{G}}$ to be the sheaf $\underline{\mathrm{Hom}}_{\mathscr{L}^{\mathrm e}}(\mathscr{L},\handwritten{\mathscr{G}})$, which is easily seen to be a subsheaf of $\handwritten{\mathscr{G}}$.

THEOREM III.7. Suppose $(\mathscr{L},\lambda)$ and $(\handwritten{\mathcal{M}},\unsure{\omega|\mathrm w})$ are two sheaves of algebras, locally separable over $(\handwritten{\mathscr{A}},\alpha)$. Then $\mathscr{L}$ is isomorphic to the commutant of $\handwritten{\mathcal{M}}$ in $\mathscr{L}\otimes_{\handwritten{\mathscr{A}}}\handwritten{\mathcal{M}}$, and $\handwritten{\mathcal{M}}$ is isomorphic to the commutant of $\mathscr{L}$.

\uline{Proof:} We will define an $\handwritten{\mathscr{A}}$-map $\mathrm{g}\colon \mathscr{L}\to \underline{\mathrm{Hom}}_{\handwritten{\mathcal{M}}^{\mathrm e}}(\handwritten{\mathcal{M}},\mathscr{L}\otimes_{\handwritten{\mathscr{A}}}\handwritten{\mathcal{M}})$.

Let $\mathrm U$ be an open set in $\mathrm X$, and $\mathrm p$ an element of $\mathscr{L}(\mathrm U)$. Define $\mathrm{g}'_{\mathrm U}(\mathrm p)$ to be the map of presheaves in $\mathrm{Hom}_{\handwritten{\mathcal{M}}^{\mathrm e\handwritten{'}}|\mathrm U}(\handwritten{\mathcal{M}}|\mathrm U,\mathscr{L}|\mathrm U\otimes_{\handwritten{\mathscr{A}}|\mathrm U}^{\handwritten{\prime}}\handwritten{\mathcal{M}}|\mathrm U)$ such that $\mathrm{g}'_{\mathrm U}(\mathrm p)_{\mathrm V}(\mathrm q)=\lambda_{\mathrm V}^{\mathrm U}(\mathrm p)\otimes\mathrm q$ for any open set $\mathrm V$ contained in $\mathrm U$ and $\mathrm q$ in $\handwritten{\mathcal{M}}(\mathrm V)$. Let $\overset{\handwritten{\rule{2.4em}{0.4pt}}}{\mathrm{g}'_{\mathrm U}(\mathrm p)}$ be the sheaf map induced by $\mathrm{g}'_{\mathrm U}(\mathrm p)$. Then $\mathrm{g}_{\mathrm U}(\mathrm p)$ will be defined as $\overset{\handwritten{\rule{2.4em}{0.4pt}}}{\mathrm{g}'_{\mathrm U}(\mathrm p)}$ for every open set $\mathrm U$ in $\mathrm X$ and $\mathrm p$ in $\mathscr{L}(\mathrm U)$.

Since $\handwritten{\mathcal{M}}$ is an $\handwritten{\mathscr{A}}$-module of finite type, it is finitely presented
% archive sheet number: 7
